\documentclass[11pt]{article}

\usepackage[margin=1in]{geometry}
\usepackage{graphicx}
\usepackage{booktabs}
\usepackage{array}
\usepackage{caption}
\usepackage{hyperref}
\usepackage{float}

\emergencystretch=3em

\graphicspath{{../all_muscle_plots/}{../top_tier_experiments/figures/}{../validation_scaling_reproducibility/figures/}{../gaitex_validation/figures/}}

\title{OpenSim-Based Estimation of Muscle-Tendon Length Trajectories\\
from Commercial IMU-Derived Gait Kinematics}
\author{\small Georgios Bouchouras\\
School of Sports Science, Department of Coaching and Physical Education\\
Metropolitan College, 46 Egnatia Str, Thessaloniki 54625, Greece\\
\textit{Corresponding author: Georgios Bouchouras}\\
\textit{Email: [INSERT EMAIL]}}
\date{}

\begin{document}

\maketitle

\noindent\textbf{Article type:} Full Paper.\par
\noindent\textbf{Main-text word count:} 2,893 words.\par
\noindent\textbf{Number of figures/tables:} 6 figures, 0 tables.

\begin{abstract}
\textbf{Background:} Wearable IMUs can assess gait outside laboratories, but they report joint kinematics rather than muscle tendon unit (MTU) behavior. We tested whether IMU joint angles can yield model based MTU length estimates comparable with optical estimates.\par
\textbf{Methods:} For primary analysis, hip, knee, and ankle angles from six trials in two participants were taken from NONAN GaitPrint, a public gait dataset of healthy older adults, and entered into OpenSim, a musculoskeletal simulation platform. Gait2392, an anatomical lower limb model mapping joint motion to muscle tendon paths, calculated ten muscle trajectories across 1397 gait cycles. Sensitivity tests examined sagittal and three dimensional inputs, 5 degree offsets, and anthropometry informed scaling. For validation, GAITEX, a public dataset with wearable inertial and optical motion capture gait recordings, provided 54 paired gait segments from 18 participants using participant specific models.\par
\textbf{Results:} All six trials produced MTU trajectories for ten muscles. Three dimensional hip motion increased gluteus medius ROM by 131\%; distal muscles changed little. Joint angle offsets changed peak and minimum lengths by less than 2\%, but ROM by up to 30.6\%. Scaling changed lengths (normalized RMSE, 70.9\% of unscaled ROM). Regeneration RMSE was below 0.001 mm. In validation, median IMU and optical correlations were 0.985 knee, 0.924 hip, and 0.884 ankle; median MTU waveform correlation was 0.880; normalized RMSE was 23.8\%.\par
\textbf{Conclusions:} IMU derived joint angles can support reproducible OpenSim MTU waveform estimation for timing and within person comparisons. They remain model derived, not direct in vivo measurements or clinical diagnostic outputs.
\end{abstract}

\noindent\textbf{Keywords:} inertial sensors; gait; OpenSim; muscle-tendon length; musculoskeletal modeling; rehabilitation

\section{Introduction}

Gait analysis is most informative when it connects what can be observed at the body surface with what is happening inside the musculoskeletal system. In the laboratory, optical motion capture and force plates provide the kinematics and kinetics required for this connection. The same laboratory pipeline, however, is difficult to repeat in a clinic, community setting, or home environment because it requires dedicated space, specialist equipment, calibration, and technical expertise.

Wearable inertial measurement units (IMUs) remove much of this logistical burden. They can record gait over long periods and can provide segment orientations or joint angle trajectories through a sensor fusion pipeline. Their practical strength is also their interpretive limitation: the usual output describes how the joints move, but not directly how a muscle-tendon unit (MTU) changes length during that movement. MTU length is therefore a potentially useful intermediate variable. It can link externally measured gait kinematics to questions about muscle stretch, shortening, limb clearance, propulsion, and compensatory movement strategies.

OpenSim provides an established way to make this link. The Gait2392 model represents lower-limb musculoskeletal geometry and maps joint coordinates to muscle-tendon paths [1,2]. Open-source inertial workflows have also shown that IMU-derived lower-limb kinematics can be obtained over prolonged gait with agreement to optical measurements [3]. Dorschky et al. integrated inertial measurements with optimal-control musculoskeletal modeling to estimate gait kinematics and kinetics [4], while OpenCap demonstrated a related portable approach based on smartphone video, model scaling, and muscle-driven simulation [5]. These studies address broader sensor-to-dynamics problems than the present work. They also make clear that a portable biomechanical result is not a single measurement: it is the product of sensor processing, kinematic reconstruction, model selection, scaling, and interpretation.

A second relevant line of work has used musculoskeletal models to estimate muscle or fiber operating lengths during gait. Arnold and Delp used experimental joint angles and activation information to estimate fiber operating lengths for 17 lower-limb muscles [6], and subsequent work examined how fiber length and velocity influence force generation across gait and running speeds [7]. These studies provide important physiological context, but they relied on laboratory-derived kinematics and carefully defined model assumptions. More recent work has shown that model choice and scaling strategy can materially affect estimated fiber lengths and excitations [8,9], while ultrasound-based studies have demonstrated that generic OpenSim geometry can differ from personalized MTU measurements [10].

This matters because a reproducible intermediate signal could help researchers and clinicians compare repeated gait sessions, identify changes in muscle lengthening and shortening, and decide when laboratory confirmation is needed. It would not replace clinical examination or direct muscle imaging, but it could make model-based hypotheses available where optical capture is impractical. The unresolved practical question is narrower and more reproducible: when joint angles are already exported by a commercial IMU system, can the subsequent conversion to MTU trajectories be made transparent, quantitatively stress-tested, and reproducible across sensor modalities? We address this question using NONAN GaitPrint, a public database of gait recordings from healthy older adults [11], and GAITEX, a public dataset that pairs wearable-sensor and optical-motion-capture recordings within the same participants [13,14]. The study does not develop raw-IMU sensor fusion, estimate muscle activation or force, or claim direct measurement of in-vivo muscle length. Instead, it evaluates a computational bridge and specifies the conditions under which its outputs should, and should not, be interpreted.

\section{Methods}

\subsection{Datasets and participants}

The primary demonstration used six four-minute gait trials from two participants in NONAN GaitPrint, a publicly available database of gait recordings from healthy older adults [11]. The public metadata identify both participants as healthy older adults in group G03. Participant S135 was a 62-year-old male, 178 cm tall and 80 kg; participant S146 was a 56-year-old male, 175 cm tall and 102 kg. The selected trials contained 288,006 frames and 1397 detected right gait cycles. These cycles represent repeated processing units rather than independent observations; the participant-level sample for the primary demonstration was n=2.

A second independent dataset, GAITEX, was used for paired modality validation. This publicly available multimodal gait dataset contains synchronized wearable inertial and optical motion capture recordings, Xsens sensor outputs, Qualisys marker trajectories, temporal gait segments, subject-specific scaled OpenSim marker models, and precomputed joint angle outputs [13,14]. Eighteen participants with complete natural gait wearable and optical recordings met the inclusion criteria for this analysis. For each participant, the first 1.5 s of each of three natural gait speed segments was analyzed, yielding 54 paired gait segments. This dataset was used to test modality agreement within participants, not to enlarge the NONAN cohort.

\subsection{Kinematic processing}

For the primary dataset, joint angles were imported as supplied by the Noraxon MyoMotion sensor fusion pipeline. No raw inertial signals were reprocessed, and no OpenSense inverse kinematics solution was generated. The primary mapping used right hip flexion, hip adduction, hip rotation, knee flexion, and ankle dorsiflexion. A sagittal-only mapping used right hip flexion, knee flexion, and ankle dorsiflexion. A secondary mapping additionally included pelvis tilt, list, and rotation. Knee and ankle frontal/transverse rotations were not included because direct equivalent coordinates were unavailable in Gait2392.

Sign conventions were checked against the direction and timing of the supplied gait-event signal. Right-foot heel strikes were used to delimit consecutive gait cycles. Each cycle was resampled to 101 equally spaced points from 0 to 100\% of the gait cycle. Time normalization was applied after segmentation, so the number of resampled points did not inflate the number of observations.

\subsection{Musculoskeletal modeling and outcomes}

MTU lengths were calculated frame by frame with the OpenSim Python API. At each sample, the selected coordinates were assigned to the model, the model was realized to position, and the muscle-tendon lengths were extracted. The primary analysis included gluteus medius, rectus femoris, vastus lateralis, tibialis anterior, soleus, and medial gastrocnemius. Four hamstrings were added to test posterior-chain extensibility: semimembranosus, semitendinosus, biceps femoris long head, and biceps femoris short head [6].

For each primary dataset trial, we calculated peak length, minimum length, ROM, and cycle-level mean and standard deviation. Repeatability was summarized by coefficients of variation across the three trials per participant. Sensitivity was quantified with waveform RMSE, Pearson correlation, and changes in ROM. Joint angle uncertainty was tested by adding uniform sagittal offsets of -5, -2, +2, and +5 degrees. Anthropometry-informed scaling used the public segment measurements and body mass for S135 and S146. Because no static calibration trial was available, this procedure was not equivalent to full marker-based subject calibration. Absolute and normalized errors were both retained because percentage errors become unstable when baseline ROM is small.

For the paired modality dataset, marker based optical inverse kinematics was regenerated from the TRC files using the marker model for each participant and lower body marker tasks. The supplied marker models in this dataset contain scaled bodies, markers, and coordinates, but no muscle actuators. We therefore transferred their lower-limb segment scale factors to a Gait2392 model that contains muscles, producing one model scaled to each participant with 92 muscle-tendon actuators per participant. Muscle paths and musculotendon parameters remained those of Gait2392. This construction personalizes skeletal segment dimensions but not muscle architecture, attachment locations, or tendon properties.

IMU and optical kinematic streams were first compared directly. Both streams were then passed through the corresponding model scaled to each participant. Agreement was summarized by Pearson correlation and RMSE normalized to the IMU-derived ROM. Median values were reported across segments because a small number of inverse kinematics segments showed large errors; complete segment level results and quality-control summaries are retained in the reproducibility bundle. An external optical-motion-capture waveform benchmark was drawn from the Dorschky dataset [12].

\section{Results}

All six NONAN trials produced smooth trajectories for the six primary muscles and four hamstrings (Figure~\ref{fig:overlays}). Regeneration of the existing primary outputs gave a maximum frame-wise RMSE of less than 0.001 mm. This result establishes reproducibility of the computational mapping from the supplied joint angle files; it does not establish measurement accuracy.

The effect of kinematic dimensionality was muscle-specific (Figure~\ref{fig:sensitivity}). Sagittal-only and three-dimensional inputs produced nearly identical trajectories for several distal muscles. Gluteus medius, by contrast, showed a mean normalized RMSE of 48.7\% and a 131.3\% mean ROM increase when hip-plane motion was included. The long hamstrings showed smaller but measurable changes. Thus, a smooth sagittal trajectory is not necessarily an adequate representation of a muscle whose geometry depends on frontal- or transverse-plane motion.

The hamstring extension produced stable waveform patterns for semimembranosus, semitendinosus, biceps femoris long head, and biceps femoris short head (Figure~\ref{fig:hamstrings}). Mean ROMs were 0.047, 0.058, 0.059, and 0.008 m, respectively. The larger hamstring excursions and late-swing peaks are consistent with their expected role in controlling limb advancement, but these model-based timings should not be interpreted as muscle activation timings. The small ROM of biceps femoris short head also explains its sensitivity to perturbation and scaling.

Uniform joint angle offsets had modest effects on peak and minimum lengths, generally below 2\%, but ROM was more sensitive (Figure~\ref{fig:uncertainty}). Maximum ROM changes reached 30.6\% for biceps femoris short head and 13.4\% for gluteus medius. Anthropometry-informed scaling preserved waveform shape while changing absolute lengths substantially; mean scaled-versus-unscaled normalized root-mean-square error (NRMSE) reached 70.9\% of unscaled ROM (Figure~\ref{fig:scaling}). These results show why generic-model waveforms may be useful for reproducibility testing but should not be read as subject-specific anatomical lengths.

The independent Dorschky normal gait benchmark gave a mean MTU waveform correlation of 0.790 and a median of 0.888 after knee sign harmonization. Because it used a different cohort, sensor system, processing pipeline, and sagittal convention, it provides external waveform plausibility rather than ground-truth validation.

The GAITEX paired modality benchmark completed for all 18 participants and 54 gait segments (Figure~\ref{fig:gaitex}). Median IMU-versus-marker correlations were 0.985 for the right knee, 0.884 for the right ankle, and 0.924 for right hip flexion. The corresponding median RMSE values were 3.36, 4.52, and 7.16 degrees. After both streams were passed through the corresponding Gait2392 models scaled to each participant, median MTU correlation across ten muscles and 54 gait segments was 0.880 and median NRMSE was 23.8\% of IMU-derived ROM. Median correlations were highest for vastus lateralis (0.977), rectus femoris (0.975), and biceps femoris long head (0.934), whereas gluteus medius showed the largest normalized disagreement. All three key right-joint RMSE values were below 10 degrees in 34 of 54 gait segments and below 20 degrees in 43 segments. In practical terms, this supports comparison of waveform shape and timing across repeated measurements, particularly for muscles with higher agreement; it does not support interpreting the plotted values as directly measured muscle lengths.

\begin{figure}[H]
\centering
\begin{minipage}{0.32\linewidth}
\includegraphics[width=\linewidth]{glut_med1_r_length_S135_vs_S146.png}
\caption*{Gluteus medius}
\end{minipage}\hfill
\begin{minipage}{0.32\linewidth}
\includegraphics[width=\linewidth]{rect_fem_r_length_S135_vs_S146.png}
\caption*{Rectus femoris}
\end{minipage}\hfill
\begin{minipage}{0.32\linewidth}
\includegraphics[width=\linewidth]{vas_lat_r_length_S135_vs_S146.png}
\caption*{Vastus lateralis}
\end{minipage}

\vspace{0.4cm}

\begin{minipage}{0.32\linewidth}
\includegraphics[width=\linewidth]{tib_ant_r_length_S135_vs_S146.png}
\caption*{Tibialis anterior}
\end{minipage}\hfill
\begin{minipage}{0.32\linewidth}
\includegraphics[width=\linewidth]{soleus_r_length_S135_vs_S146.png}
\caption*{Soleus}
\end{minipage}\hfill
\begin{minipage}{0.32\linewidth}
\includegraphics[width=\linewidth]{med_gas_r_length_S135_vs_S146.png}
\caption*{Medial gastrocnemius}
\end{minipage}
\caption{Overlay muscle-tendon length trajectories for participants S135 and S146 across the normalized gait cycle. Each panel shows one muscle, using consistent participant color coding across all six panels.}
\label{fig:overlays}
\end{figure}

\begin{figure}[H]
\centering
\includegraphics[width=0.85\linewidth]{sagittal_vs_3d_sensitivity.png}
\caption{Sensitivity of MTU waveforms to sagittal-only simplification. Bars show mean RMSE between sagittal-only and hip-3D mapping, normalized to sagittal-only ROM.}
\label{fig:sensitivity}
\end{figure}

\begin{figure}[H]
\centering
\includegraphics[width=0.90\linewidth]{hamstring_overlays.png}
\caption{Hamstring MTU length trajectories for S135 and S146 using the primary hip-3D mapping.}
\label{fig:hamstrings}
\end{figure}

\begin{figure}[H]
\centering
\includegraphics[width=0.75\linewidth]{uncertainty_rom_change.png}
\caption{Maximum absolute ROM change under uniform sagittal joint angle offsets of -5, -2, +2, and +5 degrees.}
\label{fig:uncertainty}
\end{figure}

\begin{figure}[H]
\centering
\includegraphics[width=0.85\linewidth]{scaled_vs_unscaled_sensitivity.png}
\caption{Sensitivity of MTU waveforms to anthropometry-informed scaling of Gait2392 using public NONAN subject metadata.}
\label{fig:scaling}
\end{figure}

\begin{figure}[H]
\centering
\includegraphics[width=0.94\linewidth]{gaitex_mtu_correlations.png}
\caption{Same-subject GAITEX paired modality benchmark. Panel A shows median Pearson correlation between muscle-tendon waveforms obtained from inertial and optical kinematics; panel B shows the corresponding normalized RMSE as a percentage of the inertial-sensor waveform range. Error bars indicate the interquartile range across 54 paired gait segments. Both modalities were processed with the same subject scaled Gait2392 model and downstream extraction code. Higher correlation and lower error indicate closer agreement in model-based waveform shape and amplitude.}
\label{fig:gaitex}
\end{figure}

\section{Discussion}

The study demonstrates a reproducible route from commercial IMU-derived joint angles to OpenSim MTU waveforms. Its evidentiary level is deliberately narrower than the wording often used for portable biomechanics: the pipeline is reproducible and its modality sensitivity is measurable, but the resulting lengths remain conditional model outputs rather than direct observations of muscle tissue.

The comparison with prior work clarifies both the contribution and the limitation. OpenSense provides an open-source route from IMU measurements to lower-limb kinematics and has reported agreement with optical motion capture over prolonged gait [3]. Dorschky et al. went further by embedding inertial data in an optimal-control framework for estimating gait kinematics and kinetics [4]. OpenCap addressed a different portable input, using smartphone video and a scaled musculoskeletal model to estimate dynamics and testing selected outputs against laboratory measurements [5]. The present workflow is less ambitious than those end-to-end systems: Noraxon has already performed the raw-sensor fusion, and we do not estimate forces, activations, or joint loads. Its distinct value is methodological transparency after the joint angle export. The mapping, scaling choices, perturbations, and modality comparison can be inspected without treating the commercial sensor fusion stage as if it had been developed in this study.

The GAITEX analysis strengthens the paper because it uses paired measurements from the same participants. The high knee correlation and generally strong downstream MTU correlations indicate that many waveform patterns survive the change from optical to inertial kinematics. The 23.8\% median NRMSE, however, is not negligible. Correlation primarily evaluates shape; it does not guarantee agreement in amplitude, absolute length, or phase-specific error. The result is therefore best described as preservation of model-based waveform structure under a change in sensing modality.

The model construction also requires precise language. The GAITEX skeletons were subject scaled, and their segment scale factors were retained in the Gait2392 models. This is preferable to applying an unscaled generic model to every participant. It does not make the muscles fully subject-specific. Muscle paths, attachment locations, tendon slack lengths, optimal fiber lengths, and architecture were inherited from Gait2392. The ultrasound-based work of Guenanten et al. illustrates why this distinction matters: generic and personalized MTU geometries can differ, and adding personalized insertions and paths reduced length bias [10]. Our benchmark cannot make that comparison because GAITEX does not provide independent muscle-length imaging.

The sensitivity results show why a single performance number would be misleading. Gluteus medius changed substantially when hip adduction and rotation were included, whereas several distal muscles were dominated by sagittal knee and ankle motion. A sagittal-only input can therefore be defensible for a narrowly defined distal waveform question, but it is not a general approximation for all lower-limb muscles. This interpretation is compatible with earlier model-based studies showing that joint angles and moment arms strongly influence estimated fiber operating lengths, while activation and musculotendon compliance become important when the question concerns fiber rather than whole MTU behavior [6,7].

Scaling and parameter uncertainty impose a second boundary. Falisse et al. found meaningful differences in estimated excitations and fiber lengths across model families, scaling approaches, and performance criteria, with no single combination performing best for every muscle [8]. Modenese et al. similarly showed that errors in musculotendon parameters can have muscle-specific effects during gait [9]. Our scaling analysis does not resolve these issues; it demonstrates that even segment-level personalization can materially change absolute MTU lengths. The appropriate reporting practice is consequently to separate waveform morphology, timing, absolute length, and uncertainty rather than presenting one as a substitute for the others.

The work also differs from large standardized resources such as AddBiomechanics, which provide individually scaled models and optical motion-capture/force data across hundreds of participants [15]. AddBiomechanics is valuable for population-level model benchmarking, but its muscle quantities remain model-derived. GAITEX contributes something complementary: paired IMU and optical recordings within the same participants. Neither resource supplies the full combination needed for in-vivo MTU validation, namely synchronized wearable and optical measurements together with independent muscle imaging.

The practical implication is therefore specific and bounded. The workflow can support reproducible research, exploratory rehabilitation monitoring, and hypothesis generation when the investigator is interested in relative waveform timing or within-person changes and carries model uncertainty forward. For example, it could flag a session-to-session change in late-swing hamstring lengthening or stance-phase plantar-flexor shortening for laboratory follow-up. Such a flag would be a screening signal, not a diagnosis or evidence of altered muscle structure. The workflow should not yet be used to define normative MTU ranges, diagnose pathology, prescribe treatment, or interpret absolute muscle lengths clinically. The primary NONAN demonstration includes only two participants, and the GAITEX benchmark is a separate cohort acquired with a different IMU system. A decisive next study would synchronize IMUs, optical motion capture, force plates, and ultrasound for selected muscles such as gastrocnemius medialis, then compare sensor-driven model estimates against ultrasound-tracked MTU lengths using participant-level error metrics.

\section{Conclusion}

Commercial IMU-derived joint kinematics can be converted into reproducible OpenSim MTU waveform estimates through an explicit and auditable workflow. The present evidence includes computational reproducibility, same-subject paired modality testing, hamstring extension, three-dimensional sensitivity analysis, uncertainty analysis, and construction of models scaled to each participant. It supports a model-based waveform benchmark, not in-vivo or clinical validation. The practical contribution is a way to obtain comparable muscle-tendon waveform features from repeated wearable assessments and identify cases that warrant laboratory follow-up. Future studies should add independent muscle imaging, larger cohorts, and fully personalized musculoskeletal models before absolute MTU estimates are used clinically.

\section*{Data and Code Availability}

The code supporting the workflow is available at \url{https://github.com/GiorgosBouh/opensim_project_muscle_length_imu}. A compact reproducibility bundle with scripts, public data-source information, summary outputs, and 18 muscle models scaled to each participant has been prepared for public release with the manuscript. Public data sources are NONAN (DOI: \url{10.6084/m9.figshare.27815034.v1}), Dorschky et al. (DOI: \url{10.5281/zenodo.11522050}), and GAITEX (DOI: \url{10.5281/zenodo.15729055}).

\section*{Author Contributions (CRediT)}

Georgios Bouchouras: Conceptualization, Methodology, Software, Data curation, Formal analysis, Investigation, Visualization, Writing -- original draft, Writing -- review and editing, Project administration.

\section*{Funding}

This research did not receive any specific grant from funding agencies in the public, commercial, or not-for-profit sectors.

\section*{Declaration of Competing Interest}

The author declares that there are no known competing financial or personal relationships that could have appeared to influence the work reported in this paper.

\section*{Declaration of generative AI and AI-assisted technologies in the manuscript preparation process}

During preparation of this manuscript, the author used ChatGPT to support language editing, manuscript organization, and code-oriented drafting under human oversight. The author reviewed and edited the output, verified the analyses and references, and takes full responsibility for the content of the submitted manuscript.

\section*{References}

\begin{enumerate}
\item Delp, S. L., Loan, J. P., Hoy, M. G., Zajac, F. E., Topp, E. L., and Rosen, J. M. An interactive graphics-based model of the lower extremity to study orthopaedic surgical procedures. \textit{IEEE Transactions on Biomedical Engineering} 37(8), 757--767 (1990). doi:10.1109/10.102791.
\item Seth, A. et al. OpenSim: simulating musculoskeletal dynamics and neuromuscular control to study human and animal movement. \textit{PLoS Computational Biology} 14(7), e1006223 (2018). doi:10.1371/journal.pcbi.1006223.
\item Al Borno, M. et al. OpenSense: an open-source toolbox for inertial-measurement-unit-based measurement of lower extremity kinematics over long durations. \textit{Journal of NeuroEngineering and Rehabilitation} 19, 22 (2022). doi:10.1186/s12984-022-01001-X.
\item Dorschky, E., Nitschke, M., Seifer, A. K., van den Bogert, A. J., and Eskofier, B. M. Estimation of gait kinematics and kinetics from inertial sensor data using optimal control of musculoskeletal models. \textit{Journal of Biomechanics} 95, 109278 (2019). doi:10.1016/j.jbiomech.2019.07.022.
\item Uhlrich, S. D. et al. OpenCap: Human movement dynamics from smartphone videos. \textit{PLoS Computational Biology} 19(10), e1011462 (2023). doi:10.1371/journal.pcbi.1011462.
\item Arnold, E. M. and Delp, S. L. Fiber operating lengths of human lower limb muscles during walking. \textit{Philosophical Transactions of the Royal Society B} 366, 1530--1539 (2011). doi:10.1098/rstb.2010.0345.
\item Arnold, E. M., Hamner, S. R., Seth, A., Millard, M., and Delp, S. L. How muscle fiber lengths and velocities affect muscle force generation as humans walk and run at different speeds. \textit{Journal of Experimental Biology} 216, 2150--2160 (2013). doi:10.1242/jeb.075697.
\item Falisse, A. et al. Evaluation of musculoskeletal models, scaling methods, and performance criteria for estimating muscle excitations and fiber lengths across walking speeds. \textit{Journal of Biomechanics} 143, 111289 (2022). doi:10.1016/j.jbiomech.2022.111289.
\item Modenese, L., Ceseracciu, E., Reggiani, M., and Gorton, G. E. Sensitivity of subject-specific models to Hill muscle-tendon model parameters in simulations of gait. \textit{Journal of Biomechanics} 49, 2936--2943 (2016). doi:10.1016/j.jbiomech.2016.04.008.
\item Guenanten, H. et al. Muscle-tendon unit length measurement using 3D ultrasound in passive conditions: OpenSim validation and development of personalized models. \textit{Annals of Biomedical Engineering} 52, 997--1008 (2024). doi:10.1007/s10439-023-03436-2.
\item Wiles, T. M., Kim, S. K., Mangalam, M., Sommerfeld, J. H., Brink, K. J., Grunkemeyer, A., Kalaitzi Manifrenti, M., Charles, A. E., Shakerian, N., Haghighatnejad, M., Mastorakis, S., Stergiou, N., and Likens, A. D. NONAN GaitPrint: an IMU gait database of healthy older adults. \textit{Scientific Data} 12, 143 (2025). doi:10.1038/s41597-024-04359-w.
\item Dorschky, E., Nitschke, M., Seifer, A. K., van den Bogert, A. J., and Eskofier, B. M. Lower-body inertial sensor and optical motion capture recordings of walking and running [dataset]. Zenodo (2024). doi:10.5281/zenodo.11522050.
\item Spilz, A., Oppel, H., Werner, J., Stucke-Straub, K., Capanni, F., and Munz, M. GAITEX: Human motion dataset of impaired gait and rehabilitation exercises using inertial and optical sensors. \textit{Scientific Data} 13, 11 (2025). doi:10.1038/s41597-025-06439-x.
\item Spilz, A., Oppel, H., Werner, J., Stucke-Straub, K., Capanni, F., and Munz, M. Dataset GAITEX: Human motion dataset of impaired gait and rehabilitation exercises using inertial and optical sensors [dataset]. Zenodo (2025). doi:10.5281/zenodo.15729055.
\item Werling, K. et al. AddBiomechanics Dataset: Capturing the Physics of Human Motion at Scale. In: \textit{Computer Vision -- ECCV 2024}, Lecture Notes in Computer Science 15146, 490--508 (2025). doi:10.1007/978-3-031-73223-2\_27.
\end{enumerate}

\end{document}
